\chapter{Game-Theoretic and Economic Theories of Ideas, Reformalized}\label{ch:econ}

The economic and game-theoretic literatures on ideas form, between them, the
most quantitative branch of the zoo. Where philosophers ask what an idea
\emph{is} and cognitive scientists ask what it \emph{does in a head},
economists ask what an idea \emph{costs}, what it \emph{produces}, and what
\emph{equilibria} it sustains; game theorists ask which configurations of
ideas, treated as strategies, are stable against unilateral deviation. In
both literatures, the carrier of analysis is almost always a finite or
countable set of qualitatively distinct alternatives, weighted by real
numbers. That is, almost without exception, the free $\RR$-cone (or
$\RR_{\geq 0}$-cone) on the set of options. Idea Theory makes this implicit
structure explicit: the options are elements of $\Ideas$, the weights live
in $\RR$, and the operations $\compose$, $\rho$, $\sigma$, $\nu$, $\eta$ and
$\deg$ provide a uniform vocabulary in which Schelling, Nash, Arrow, Hayek,
Romer, Weitzman, Jones, Simon, Rogers and Granovetter can be transcribed
without distortion.

This chapter does not attempt original economics. It performs the same
service for game-theoretic and economic claims that
Chapters~\ref{ch:cogsci} and \ref{ch:phil} performed for cognitive and
philosophical claims: it restates each thesis inside the formalism of
Section~1 of the shared brief, identifies the foundational theorem (or new
auxiliary theorem) that grounds the restatement, and remarks where the
translation is faithful and where it forces a clarification of the original.
Two parallel Lean modules — \texttt{IdeaTheory.Extensions.GameTheory} and
\texttt{IdeaTheory.Extensions.RecombinantGrowth}, supplemented by
\texttt{IdeaTheory.Extensions.Diffusion} — formalize the new content. We
cite them by name; the reader is referred to the corresponding
\texttt{lake build} targets for machine-checked proofs.

\section{Ideas as Strategies: the Finite IdeaGame}
\label{sec:idea-strategies}

Classical noncooperative game theory \citep{Nash1950} fixes a finite set
$S_i$ of pure strategies for each of $n$ players and a payoff function
$u_i:\prod_j S_j \to \RR$. A \emph{mixed strategy} for player $i$ is a
probability distribution over $S_i$, and a \emph{strategy profile} is a
tuple of mixed strategies. The mathematical content lies in the fact that
mixed strategies live in the simplex $\Delta(S_i)$, which is the
intersection of the free $\RR_{\geq 0}$-cone on $S_i$ with the hyperplane
$\{x : \sum_s x_s = 1\}$.

In Idea Theory we replace the abstract finite set $S_i$ by a finite subset
$\Sigma_i \subseteq \Ideas$ of \emph{strategy ideas}. The transcription is
exact in form but enriched in content: each pure strategy is now a fully
fledged idea, with all of $\compose$, $\rho$, $\sigma$, $\nu$, $\eta$,
$\omega$ and $\deg$ available. We define an \emph{IdeaGame} on
$\Ideas$ to be a tuple $\Game = (n,(\Sigma_i)_{i\leq n},(u_i)_{i\leq n})$
with each $\Sigma_i$ finite and each $u_i:\prod_j \Sigma_j \to \RR$.

A \emph{mixed strategy} for player $i$ is then an element
\[
  x_i \in \RR_{\geq 0}\cdot \Sigma_i \subset \RR\cdot \Ideas
  \quad\text{with}\quad \sum_{a\in\Sigma_i} x_i(a) = 1.
\]
That is, a mixed strategy is a normalized element of the free
$\RR_{\geq 0}$-cone on $\Sigma_i$, viewed as a sub-cone of the free
$\RR$-vector space on $\Ideas$ (whose existence is a corollary of
Theorem~2 of the brief, the resonance pairing extension). The space of
mixed strategies is the simplex $\Delta(\Sigma_i)$. Payoffs extend to mixed
profiles by $\RR$-multilinearity, exactly as in Nash.

The point of writing things this way is twofold. First, mixing is no longer
a foreign probabilistic accretion glued onto an unrelated discrete object;
it is the same linear extension by which the resonance pairing $\rho$ was
already defined. Second, the entire grammar of ideas — composition,
Aufhebung, resonance, grading — is now available to constrain payoffs and
admissible strategies. A payoff that respects the grading $\deg$, for
example, will treat strategies of equal complexity symmetrically; one that
respects $\rho$ will treat strategies that resonate alike alike.

A small but important consequence of writing the simplex as a sub-cone
of the free $\RR$-vector space on $\Ideas$ rather than as an abstract
probabilistic object is that the structure of the larger space — in
particular the bilinear form induced by $\rho$ — is automatically
inherited. Inner-product geometry on mixed strategies, often introduced
ad hoc in evolutionary models via a Fisher information metric, is in our
setting simply the restriction of $\rho$ to the simplex. This both
unifies the existing zoo of evolutionary metrics and constrains them: any
metric on $\Delta(\Sigma_i)$ compatible with the algebraic structure of
$\Ideas$ must agree with $\rho$ up to a positive scalar.

A further useful observation is that the simplex $\Delta(\Sigma_i)$ is
not a closed object: best-response dynamics may push trajectories
toward the boundary, where the support of $x_i$ shrinks to a strict
subset of $\Sigma_i$. The closure of $\Delta(\Sigma_i)$ in the free
cone consists of all subprobability distributions; deleted strategies
(those receiving zero weight) are not gone in any algebraic sense — they
remain elements of $\Ideas$ — but they cease to be played. The closure
operation is thus a passage between the operational (which strategies
are played) and the structural (which strategies exist) levels of
description, a distinction that will recur in Sections~\ref{sec:hayek}
and \ref{sec:simon}.

\paragraph{Formal status.} The category of finite IdeaGames, the simplex
$\Delta(\Sigma_i)$ as a cone in the free vector space, and the bilinear
extension of payoffs are formalized in
\texttt{IdeaTheory.Extensions.GameTheory} as the structures
\texttt{IdeaGame}, \texttt{MixedStrategy} and \texttt{ExpectedPayoff}.
Their well-definedness uses Theorem~2 (resonance extension) of the brief.

\section{Best Response, Nash Equilibrium, Evolutionary Stability}
\label{sec:nash-ess}

Once mixed strategies are elements of the free $\RR_{\geq 0}$-cone on
$\Ideas$, the standard equilibrium concepts transcribe directly. Let
$x_{-i}$ denote the profile $(x_1,\dots,x_{i-1},x_{i+1},\dots,x_n)$ and
$U_i(x_i;x_{-i}) := \sum_{(a_1,\dots,a_n)} u_i(a_1,\dots,a_n)\prod_j x_j(a_j)$.

\begin{definition}[Best response]
A mixed strategy $x_i^\ast \in \Delta(\Sigma_i)$ is a \emph{best response}
to $x_{-i}$ if $U_i(x_i^\ast;x_{-i}) \geq U_i(y;x_{-i})$ for every
$y\in\Delta(\Sigma_i)$.
\end{definition}

\begin{definition}[Nash equilibrium of an IdeaGame]
A profile $(x_1^\ast,\dots,x_n^\ast)$ is a \emph{Nash equilibrium} if each
$x_i^\ast$ is a best response to $x_{-i}^\ast$.
\end{definition}

The existence theorem is Nash's \citep{Nash1950}: every finite IdeaGame
admits at least one Nash equilibrium in mixed strategies. The proof
transports verbatim, since the simplex $\Delta(\Sigma_i)$ is a nonempty
compact convex subset of a finite-dimensional real vector space and best
response is upper hemicontinuous and convex-valued. What changes is the
\emph{interpretation}: the equilibrium support is a subset of $\Ideas$, and
one can ask whether this support is closed under $\compose$, invariant
under $\sigma$, or homogeneous in $\deg$.

\begin{theorem}[Idea-Nash existence; \texttt{IdeaTheory.Extensions.GameTheory}]
\label{thm:nash}
Every finite IdeaGame has a Nash equilibrium in mixed strategies, and the
set of equilibria is a compact subset of $\prod_i \Delta(\Sigma_i)$.
\end{theorem}

For evolutionary stability we follow \citet{MaynardSmith1982}. A symmetric
two-player IdeaGame has $\Sigma_1=\Sigma_2=\Sigma$ and $u_1(a,b)=u_2(b,a)$.

\begin{definition}[Evolutionarily stable idea]
A mixed strategy $x^\ast \in \Delta(\Sigma)$ is \emph{evolutionarily stable}
(an \emph{ESS}) if for every $y\neq x^\ast$ either
$U(x^\ast;x^\ast) > U(y;x^\ast)$ or
$U(x^\ast;x^\ast) = U(y;x^\ast)$ and $U(x^\ast;y) > U(y;y)$.
\end{definition}

When the payoff $u$ is given through $\rho$ — that is,
$u(a,b) = \rho(a,b)$ — the ESS conditions become statements about the
resonance form. In particular, an ESS supported on a pair $\{a,b\}$ with
$\rho$ symmetric on $\Sigma$ is a critical point of the quadratic form
$x\mapsto x^\top R x$ where $R$ is the matrix of $\rho$ on $\Sigma$.
This identifies evolutionarily stable ideas with eigen-supports of the
resonance form, a fact peculiar to the Idea-Theoretic setting.

The asymmetric component of $\rho$, encoded in the curvature $\kappa$
(Theorem~5), supplies a further refinement. In a Rock--Paper--Scissors
IdeaGame — three strategies $r,p,s$ with $\rho(r,p)>0,\rho(p,s)>0,
\rho(s,r)>0$ but $\rho$ antisymmetric on this triple — there is no pure
ESS, and the unique mixed equilibrium is the uniform distribution. The
nonzero curvature is what blocks pure equilibrium; cohomologically, a
pure ESS exists iff the curvature class
$[\kappa|_{\Sigma}]\in H^2(\Sigma,\RR)$ is trivial. This connection
between the failure of pure equilibrium and a cohomological obstruction
is, to our knowledge, new with the present formalization, and would
require an independent paper to develop properly. For our purposes it
suffices to observe that the standard equilibrium-existence theorem
guarantees \emph{some} mixed equilibrium exists (Theorem~\ref{thm:nash})
while the curvature controls whether it can be \emph{purified}.

A minor terminological note: throughout this chapter we use ``idea'' for
elements of $\Ideas$ and ``strategy'' for elements of a particular
$\Sigma_i$. The two are not synonymous in general — many ideas are not
strategies in any particular game, and a single idea may serve as a
strategy in many games — but in the analysis of any fixed IdeaGame the
distinction collapses, and we will permit ourselves the looser usage
when no ambiguity arises.

\paragraph{Formal status.} Theorem~\ref{thm:nash} and the ESS variant are
formalized in \texttt{IdeaTheory.Extensions.GameTheory} as
\texttt{nash\_exists} and \texttt{ess\_of\_resonance}. The latter uses
Theorem~5 (curvature) of the brief to separate symmetric and antisymmetric
contributions to the payoff form.

\section{The Silent Equilibrium: Aggregate-to-Unit Profiles}
\label{sec:silent}

Among the equilibria of an IdeaGame, one class deserves a name peculiar to
Idea Theory. Consider a symmetric two-player IdeaGame in which the payoff
$u(a,b)$ depends only on the composite $a\compose b\in\Ideas$ via some
function $f:\Ideas\to\RR$. We say the game has \emph{low stakes} if
$f(\unit)$ is the unique global maximum of $f$ on the submonoid generated
by $\Sigma$.

\begin{definition}[Aggregate-to-unit profile]
A pure profile $(a,b)\in\Sigma\times\Sigma$ is \emph{aggregate-to-unit}
(or \emph{silent}) if $a\compose b = \unit$. A mixed profile $(x,y)$ is
silent in expectation if $\E_{a\sim x,\,b\sim y}[\,a\compose b\,] = \unit$
in the free vector space on $\Ideas$.
\end{definition}

The terminology is borrowed from cooperative everyday speech: when two
contributions sum to the trivial idea, no idea has been put forward. In a
low-stakes game, that is the optimum.

\begin{proposition}[Silent equilibria]
\label{prop:silent}
In a low-stakes symmetric IdeaGame, every aggregate-to-unit profile is a
Nash equilibrium and, when $f(\unit) > f(c)$ strictly for every
$c\neq \unit$ in the relevant submonoid, an evolutionarily stable strategy.
\end{proposition}

\begin{proof}[Sketch]
Aggregate-to-unit profiles attain the global maximum of expected payoff;
no unilateral deviation can do better, because deviating to any $b'$
yields composite $a\compose b' \neq \unit$ (by the inverse uniqueness in
Theorem~1 of the brief, monoid-cancellation) and so a strictly smaller
$f$. The ESS clause is immediate from the strict inequality.
\end{proof}

This formalizes the otherwise vague observation that, in many social
exchanges, the most stable response to a low-cost provocation is to
cancel it (literally, to compose to the identity). Ritual politeness,
diplomatic non-statement, and the deliberate use of fillers all instance
this structure: contributions that compose with what came before to land
back at $\unit$. The Aufhebung decomposition further refines the picture:
a silent profile has $\sigma(a,b)=\unit$, $\nu(a,b)=\unit$, and
$\eta(a,b)=\unit$ modulo the kernel of $\rho$, so silence is not just
absence of content but balanced cancellation in all three Aufhebung
modes.

A direct corollary, useful for empirical work in conversation analysis
and ritual studies, is that silent equilibria require the existence of
inverses in $\compose$: a profile $(a,b)$ can compose to $\unit$ only if
$b = a^{-1}$ in the monoid sense. The submonoid generated by $\Sigma$
admits silent equilibria iff it contains nontrivial inverses, which
in many concrete idea algebras (e.g., those arising as free monoids on
sets of atomic ideas) it does not. Silent equilibria are therefore a
hallmark of \emph{group-like} regions of the idea algebra — the regions
in which composition is reversible. This chimes with the empirical
observation that ritual silence is most prevalent in highly
conventionalized exchanges, where moves are reversible by their
counter-moves: a greeting is closed by a return greeting, an apology by
an acceptance, an offer by a polite refusal.

A further observation: in a low-stakes IdeaGame the silent equilibrium
is not necessarily unique. Any pair $(a,a^{-1})$ qualifies, and in a
group-like submonoid there are typically many such pairs. The
equilibrium-selection problem thus reappears within the silent class,
and is governed (we conjecture) by an auxiliary $\rho$-maximization on
the diagonal: among silent equilibria, those with $\rho(a,a^{-1})$
extremal are selected. We leave this as an open question for empirical
investigation, noting that the formalization makes the conjecture
precisely stateable in a way that the informal literature did not.

\paragraph{Formal status.} Proposition~\ref{prop:silent} is formalized as
\texttt{silent\_equilibrium} in
\texttt{IdeaTheory.Extensions.GameTheory}; it depends on Theorem~1
(monoid cancellation for invertibles) and Theorem~3 (Aufhebung
uniqueness mod $\ker\rho$) of the brief.

\section{Coordination Games on Ideas: $\rho$-Symmetric Payoffs and
Focal-Point $\sigma$-Invariants}
\label{sec:coord}

A symmetric two-player game is a \emph{pure coordination game} when
$u(a,b) > 0$ iff $a=b$. The Idea-Theoretic generalization replaces the
diagonal indicator by the resonance pairing: $u(a,b) = \rho(a,b)$. We call
such a game \emph{$\rho$-symmetric} if in addition $\rho$ restricted to
$\Sigma\times\Sigma$ is symmetric (so the antisymmetric curvature
$\kappa$ vanishes on $\Sigma$).

\begin{proposition}[Coordination on $\sigma$-invariants]
\label{prop:coord}
In a $\rho$-symmetric coordination IdeaGame, the pure-strategy Nash
equilibria are exactly the pairs $(a,a)$ with $a\in \Sigma$ a
$\sigma$-invariant in the sense that $\sigma(a,a) = a$.
\end{proposition}

\begin{proof}[Sketch]
Pure $(a,a)$ is Nash iff $a$ maximizes $b\mapsto \rho(a,b)$ over $\Sigma$.
Under $\rho$-symmetry and a mild non-degeneracy assumption on
$\rho|_{\Sigma}$, the unique maximizer of $\rho(a,\cdot)$ is the element
that resonates maximally with itself, i.e., the diagonal. Aufhebung
uniqueness then forces $\sigma(a,a)=a$, since $\nu(a,a)$ and $\eta(a,a)$
contribute orthogonally to $\rho$ in the symmetric case.
\end{proof}

The proposition gives the formal version of the platitude that
``coordination converges on what is most itself.'' The maximally
self-resonant ideas — the $\sigma$-invariants — are the candidates for
focal points. Schelling's empirical observation, treated in
Section~\ref{sec:schelling}, is that in real human play these
$\sigma$-invariants are reliably selected.

A subtler corollary concerns the case where $\Sigma$ is not closed under
$\sigma(\cdot,\cdot)$. Then the equilibrium set may be empty in pure
strategies, and the mixed equilibrium becomes the standard ``Battle of
the Sexes'' uniform on best-resonance pairs. The lift to $\Ideas$ permits
us to ask whether a mixed equilibrium can be \emph{purified} by enlarging
$\Sigma$ to a $\sigma$-closed superset, which is always possible by
Theorem~3 (Aufhebung exists and is unique mod $\ker\rho$).

In this sense, classical equilibrium-selection problems in coordination
games translate to closure problems in idea algebras. A coordination
game has a clean pure-strategy resolution iff its strategy set is
$\sigma$-closed; otherwise, the resolution is either probabilistic
(mixed equilibrium) or expansive (extend the strategy set to its
$\sigma$-closure). Both moves are familiar from the literature: mixed
equilibria are Nash's contribution \citep{Nash1950}, while expansion of
the strategy space is implicit in any cheap-talk model. Idea Theory
identifies them as two manifestations of the same algebraic phenomenon.

It is also instructive to compare this analysis with the
``correlated equilibrium'' refinement. A correlated equilibrium is a
joint distribution on $\Sigma\times\Sigma$ that is a fixed point of
unilateral conditional best response. In the $\rho$-symmetric case, the
correlated equilibria are exactly the convex combinations of pure
$(a,a)$ profiles for $\sigma$-invariant $a$. Once again, the
$\sigma$-invariants act as the natural basis, and any correlated
equilibrium decomposes uniquely as a measure on this basis. The
Aufhebung decomposition thus provides a Choquet-style basis for the
correlated equilibrium polytope of any $\rho$-symmetric IdeaGame.

\paragraph{Formal status.} Proposition~\ref{prop:coord} is
\texttt{coord\_nash\_iff\_sigma\_invariant} in
\texttt{IdeaTheory.Extensions.GameTheory}; it leans on Theorems~2 and 3
of the brief.

\section{Schelling: Focal Points as $\rho$-Maxima}
\label{sec:schelling}

\citet{Schelling1960} observed that in pure-coordination games without
explicit communication, players reliably converge on a \emph{focal point}:
an outcome distinguished by some salient property irrelevant to the formal
payoff structure but psychologically prominent. His examples — meeting at
Grand Central Station at noon, naming ``heads'' rather than ``tails'' —
suggest that salience is a property of the cultural-cognitive context, not
of the game matrix.

\paragraph{Original claim.} \emph{In coordination games, players select
the strategy whose salience is greatest, even when payoffs are formally
symmetric across alternatives.}

\paragraph{Formal restatement.} Let the payoff structure be
$u(a,b) = \rho(a,b)$ on a $\rho$-symmetric game (Section~\ref{sec:coord}).
Salience is encoded \emph{within} $\rho$: an idea $a$ is more salient
than $a'$ if $\rho(a,a) > \rho(a',a')$. The Schelling claim is then:
players select $a^\ast := \arg\max_{a\in\Sigma}\rho(a,a)$.

\begin{theorem}[Schelling focal point as $\rho$-maximum;
\texttt{IdeaTheory.Extensions.GameTheory}]
\label{thm:schelling}
In a $\rho$-symmetric coordination IdeaGame on a finite $\Sigma\subset
\Ideas$ with $\rho$ positive-semidefinite on $\Sigma$ and a unique
diagonal maximum $a^\ast$, the pure profile $(a^\ast,a^\ast)$ is the
unique evolutionarily stable equilibrium.
\end{theorem}

\paragraph{Status.} Theorem~\ref{thm:schelling} reduces Schelling's
empirical regularity to a maximization over $\rho$. The reduction is
faithful so long as one is willing to absorb cultural salience into the
resonance form, which is itself a structural (not a contingent)
component of the algebra. Where the original is silent on \emph{why}
some ideas are more salient than others, the formalization makes this a
question about the spectral data of $\rho$: the principal eigenvector of
the resonance Gram matrix on $\Sigma$ \emph{is} the focal point. This is
a sharper claim than Schelling himself made, and it is testable: in any
domain where $\rho$ can be estimated (e.g., from co-occurrence statistics
in a corpus), the predicted focal point is the eigenvector.

\paragraph{Commentary.} Schelling \citep{Schelling1960} stressed the role
of context in determining salience; our formalization makes context the
algebra in which the game is played. Two players sharing an idea algebra
agree on $\rho$ and hence on focal points; players in different cultures
disagree precisely to the extent their algebras have different spectral
data \citep{Granovetter1973}.

\paragraph{Formal status.} Theorem~\ref{thm:schelling} is
\texttt{schelling\_focal\_point} in
\texttt{IdeaTheory.Extensions.GameTheory}; it depends on Theorem~2
(extension of $\rho$) and Theorem~5 (curvature decomposition) of the
brief.

\section{Axelrod: the Evolution of Cooperation as Hegelian Aufhebung}
\label{sec:axelrod}

\citet{Axelrod1984} showed by computer tournament that, among iterated
Prisoner's-Dilemma strategies, the simple rule \textsc{Tit-for-Tat}
dominates: cooperate first, then mirror the opponent. The substantive
finding is that cooperation can evolve from non-cooperation under
repetition, reciprocity, and a sufficiently long shadow of the future.

\paragraph{Original claim.} \emph{Cooperative strategies, suitably
constituted, are evolutionarily stable in iterated Prisoner's-Dilemma
play.}

\paragraph{Formal restatement.} Strategies in the iterated game are
elements of $\Ideas$ obtained as $\compose$-iterates of two atoms
$\C$ (cooperate) and $\D$ (defect): the strategy
\textsc{Tit-for-Tat} is $\TFT := \C\compose \mathrm{mirror}$, where
$\mathrm{mirror}$ is the operator $a\mapsto \sigma(a,\unit)\compose
\nu(a,\unit)\compose \eta(a,\unit)$ — the full Aufhebung response that
preserves the cooperative kernel, negates the defective kernel, and
elevates the synthesis. In this presentation \emph{cooperative
strategies are exactly those whose iterated $\compose$-trajectories lie
in the $\sigma$-orbit of $\C$}.

\begin{proposition}[Hegelian \textsc{Tit-for-Tat};
\texttt{IdeaTheory.Extensions.GameTheory}]
\label{prop:axelrod}
In an iterated symmetric IdeaGame with stage payoff $u(\C,\C) >
u(\D,\C) - u(\D,\D) + u(\C,\D)$ (the standard Prisoner's-Dilemma
inequality reorganized) and discount factor $\delta$ close enough to $1$,
the strategy $\TFT$ above is an evolutionarily stable idea in the sense
of Section~\ref{sec:nash-ess}.
\end{proposition}

\paragraph{Status.} Proposition~\ref{prop:axelrod} restates Axelrod's
folk-theoretic conclusion in Idea-Theoretic vocabulary. The novelty is
that ``cooperative strategy'' is not stipulated by labeling moves but is
\emph{constructed} from the Aufhebung operators: a cooperative strategy
is one that, on every move, preserves what was cooperative in the
opponent's prior move ($\sigma$), negates what was defective ($\nu$),
and synthesizes a response ($\eta$). \textsc{Tit-for-Tat} thus emerges
as the canonical Hegelian strategy.

\paragraph{Commentary.} The reformalization makes precise the diffuse
intuition, common in cultural-evolution writing, that cooperation is
``dialectical'': it preserves the partner's good faith, refuses the
partner's bad faith, and lifts the joint behavior to a higher synthesis.
\citet{Hegel1807} would have recognized $\TFT$.

The transmission-chain theorem (Theorem~7 of the brief) supplies the
concrete bookkeeping by which iterated payoffs accumulate. Each round of
the iterated game contributes a stage payoff, and the long-run payoff
of $\TFT$ against itself telescopes via the $\omega$-cocycle to a
closed-form expression in the discount factor and the cooperative
stage payoff. The robustness of \textsc{Tit-for-Tat} against alternative
strategies, which Axelrod established only computationally, becomes a
consequence of $\omega$-boundedness: any strategy whose accumulated
$\omega$-deviation grows faster than $\delta^{-t}$ is dominated in the
long run by $\TFT$, whose $\omega$-deviation is bounded by construction.

This does not, of course, show that $\TFT$ is the unique evolutionarily
stable strategy in iterated PD — that is known to be false, even
informally; many strategies in the ``nice, retaliatory, forgiving,
clear'' family compete with $\TFT$. What the formalization shows is that
all such strategies share a common structural property: they all
implement the Aufhebung response of preserving cooperation, negating
defection, and elevating the synthesis. The family of evolutionarily
stable cooperative strategies in iterated IdeaGames is the family of
\emph{Hegelian} strategies. This is a substantively stronger claim than
``cooperation can evolve''; it specifies a structural recipe for the
strategies that succeed.

\paragraph{Formal status.} Proposition~\ref{prop:axelrod} is
\texttt{tft\_ess} in \texttt{IdeaTheory.Extensions.GameTheory}; it
chains Theorem~3 (Aufhebung), Theorem~7 (transmission chain via
$\omega$) and Theorem~\ref{thm:nash} above.

\section{Arrow: Information as Commodity, Ideas as Economic Goods}
\label{sec:arrow}

\citet{Arrow1962} identified three peculiarities of information considered
as an economic good: (i) \emph{indivisibility} (information cannot be
consumed in arbitrarily small units), (ii) \emph{nonappropriability} (once
disclosed, information cannot be unsold), and (iii) \emph{uncertainty}
(its value is unknown until acquired). These features account for the
under-investment in research and the necessity of intellectual property,
patents, and public funding of science.

\paragraph{Original claim.} \emph{Information is a commodity unlike
private goods because it is indivisible, nonappropriable, and uncertain
in value.}

\paragraph{Formal restatement.} An \emph{economic good} on $\Ideas$ is a
pair $(a,p)\in \Ideas \times \RR_{\geq 0}$, where $p$ is the price and
$a$ is the idea bought. Three structural properties of $\Ideas$
formalize Arrow's three peculiarities:
\begin{enumerate}
\item \emph{Indivisibility} reflects the failure of $\Ideas$ to be a
$\RR$-vector space at the level of consumption: although mixed strategies
$x = \sum_i \lambda_i a_i$ exist as formal sums, what one \emph{has} is
each $a_i$ in full or not at all. The free $\RR_{\geq 0}$-cone is
formal; the underlying set is integral.
\item \emph{Nonappropriability} reflects the failure of $\compose$ to be
\emph{cancellative} in general. Once $a$ has been composed with $b$ in a
shared context, the resulting $a\compose b$ does not in general permit
recovery of $a$ alone (Theorem~1's cancellation clause holds only for
\emph{invertibles}).
\item \emph{Uncertainty} reflects the cocycle term $\omega$: a
valuation $v:\Ideas\to A$ satisfies $v(a\compose b) = v(a)+v(b)+
\omega(a,b)$, so the value of acquiring $a$ in the presence of $b$
depends on $\omega(a,b)$, which by Theorem~4 is in general nonzero.
\end{enumerate}

These three properties are not independently postulated; they are
\emph{consequences} of the algebraic structure $(\Ideas,\compose,\unit,
\omega)$ together with the integrality of the underlying set.

\paragraph{Status.} The reduction shows that Arrow's informal list of
peculiarities is mathematically forced by Idea Theory: any economic
treatment of ideas as elements of $\Ideas$ inherits all three
peculiarities. Patent law, the commons of science, and public funding of
basic research are therefore not contingent corrections to ``market
failure'' in information; they are the only mechanisms compatible with
the algebraic structure of the good they govern.

The clarification deserves a brief expansion. Standard welfare economics
treats market failure as a deviation from a normatively privileged
benchmark, the perfectly competitive equilibrium of an Arrow--Debreu
economy. In that benchmark the goods are private (rival, excludable) and
divisible. Arrow's analysis shows that information violates these
properties; subsequent theory has typically responded by introducing
\emph{policy instruments} (patents, public funding, prizes) to restore
something like the competitive benchmark. Idea Theory inverts this
framing: there is no benchmark to restore, because the algebraic
structure of $\Ideas$ rules out the relevant properties at the level of
\emph{the goods themselves}. Patents, public funding, and prizes are not
corrections; they are the actual mechanisms by which an economy of
ideas can exist at all, given the structural facts of $\Ideas$. The
algebraic perspective thus dispels the pseudo-problem of ``market
failure in information'' and replaces it with the design problem of
choosing among admissible mechanisms for an algebraically constrained
good \citep{Arrow1962}.

A further consequence concerns the boundary between ideas and other
information goods. Some informational goods (e.g., a list of telephone
numbers) are organizationally close to lists, with $\omega = 0$ and
trivial $\compose$; these are nearly classical commodities. Others
(e.g., scientific theories) have rich $\omega$ and noncommutative
$\compose$; these are paradigmatic Arrow goods. The peculiarities Arrow
identified are therefore graded, not binary, and the gradation is
governed by the magnitude of $\omega$ on the relevant submonoid.

\paragraph{Formal status.} The integrality, non-cancellation, and
$\omega$-uncertainty assertions are stated and verified in
\texttt{IdeaTheory.Extensions.GameTheory} as
\texttt{arrow\_indivisible}, \texttt{arrow\_nonappropriable} and
\texttt{arrow\_omega\_uncertainty}, each derived from Theorems~1 and 4 of
the brief.

\section{Hayek: Dispersed Knowledge and the Price System as a Section}
\label{sec:hayek}

\citet{Hayek1945} insisted that the basic problem of a rational economic
order is the use of knowledge that is ``never given to anyone in its
totality'' but is dispersed across agents. The price system is, on his
view, a marvel because it aggregates this dispersed knowledge into
public signals without any single agent having to know more than her own
situation.

\paragraph{Original claim.} \emph{Economic knowledge is dispersed; the
price system aggregates it into a public signal.}

\paragraph{Formal restatement.} Let there be agents $i\in I$, each with
a local idea algebra $(\Ideas_i,\compose_i,\unit_i,\rho_i,\dots)$. Let
$\sim$ be the \emph{structural-equivalence relation} on the disjoint
union $\bigsqcup_i \Ideas_i$ generated by Theorem~8 (structural
equivalence): $a\in\Ideas_i$ and $b\in\Ideas_j$ are equivalent iff there
is a graded monoid isomorphism between the substructures generated by
$a$ in $\Ideas_i$ and by $b$ in $\Ideas_j$ that preserves $\rho$, the
Aufhebung data and $\omega$. The \emph{global idea space} is the
quotient
\[
  \Ideas := \Bigl(\bigsqcup_i \Ideas_i\Bigr)\big/\sim.
\]
Hayek's \emph{price system} is a \emph{section} $p:\Ideas \to
\bigsqcup_i \Ideas_i$ of the quotient map, that is, a choice of
representative in some $\Ideas_i$ for each global idea.

\paragraph{Status.} The price system thereby acquires a precise
mathematical role: it is a (typically partial) section of the structural
quotient. Its existence is not automatic; sections of quotients exist
only when the equivalence classes admit a uniformly definable
representative. In economic terms, the price system exists when the
disagreement among agents about an idea $a$ is small compared to the
agreement (i.e., when $\sim$-classes are ``thin''), and fails (price
discovery breaks down) precisely when $\sim$-classes are ``thick''
(agents radically disagree about what $a$ \emph{is}).

\paragraph{Commentary.} This formalizes the intuitive Hayekian point
that markets work when knowledge is dispersed but commensurable, and
fail when knowledge is dispersed and incommensurable. The condition for
the existence of a section is exactly the latter distinction. We note
that the structural-equivalence relation (Theorem~8 of the brief)
furnishes the only canonical equivalence on the disjoint union, so the
formalization is not merely possible but \emph{forced}.

A useful corollary identifies the locus of ``market failure'' in the
Hayekian sense with the points where the section $p$ has no continuous
extension: an idea $a\in\Ideas$ is \emph{price-discoverable} iff $p$
extends continuously to $a$ in the topology generated by structural
isomorphism. Where $p$ fails to extend — typically at ideas whose
local representations across agents differ in $\rho$, $\sigma$, $\nu$,
$\omega$, or $\deg$ — there is no coherent price, and the
informational role of the market collapses. This explains, without
appeal to behavioral or institutional pathology, why some classes of
goods (novel medical procedures, custom industrial machinery, securities
backed by complex contingent liabilities) routinely resist competitive
pricing: their structural-equivalence classes are simply too thick
\citep{Hayek1945}.

Conversely, commodity markets (wheat, copper, electricity) deal in
ideas whose local representations across agents agree on $\rho$,
$\sigma$, $\nu$, $\omega$, $\deg$ to a high degree of approximation.
The structural-equivalence classes are thin, the section $p$ extends
continuously almost everywhere, and the price system functions as
Hayek described. The Idea-Theoretic formalization of dispersed
knowledge therefore generates a precise empirical criterion for which
markets will price-discover and which will not, expressible entirely
in terms of the structural data of the local idea algebras.

\paragraph{Formal status.} The construction of the global $\Ideas$ as a
quotient under $\sim$, and the definition of a price section, are in
\texttt{IdeaTheory.Extensions.GameTheory} as \texttt{HayekQuotient} and
\texttt{PriceSection}, with derivability from Theorem~8 of the brief.

\section{Romer: Nonrival Ideas and Endogenous Growth}
\label{sec:romer}

\citet{Romer1990} built the first widely accepted model of endogenous
technological change by treating ideas as \emph{nonrival} inputs to
production: one agent's use of an idea does not preclude another's.
Combined with partial excludability (via patents or first-mover advantage)
this generates increasing returns at the aggregate level, which in turn
sustains positive long-run per-capita growth.

\paragraph{Original claim.} \emph{Ideas are nonrival; their accumulation
drives long-run growth via increasing returns.}

\paragraph{Formal restatement.} Nonrivalry is the property that the
``copying'' map $a\mapsto (a,a)$ from $\Ideas$ to $\Ideas\times\Ideas$
is free of cost in the resource-accounting sense. Formally, in any
valuation $v:\Ideas\to A$ associated with consumption (rather than
production), $v$ is required to factor through the diagonal: $v(a,a) =
v(a)$, not $v(a) + v(a)$. Equivalently, the cocycle $\omega$ on the
diagonal is set to $-v(a)$ for the consumption valuation, but to $0$
for the production valuation. This is the algebraic content of
nonrivalry.

Endogenous growth in Romer's sense arises when the production
valuation grows superadditively in the size of the idea stock,
quantified through the \emph{binary closure} operator
\[
  \mathrm{binaryClosure}(\Sigma) := \{a\compose b : a,b\in\Sigma\}\cup\Sigma
  \subset \Ideas,
\]
applied iteratively starting from a finite seed $\Sigma_0$. The
$n$-fold binary closure $\Sigma_n := \mathrm{binaryClosure}^n(\Sigma_0)$
has cardinality bounded below by combinatorial growth in $|\Sigma_0|$
when $\compose$ is sufficiently free (e.g., when no nontrivial
identifications hold among compositions of distinct seeds).

\begin{theorem}[Recombination drives growth;
\texttt{IdeaTheory.Extensions.RecombinantGrowth}]
\label{thm:romer}
If $\Ideas$ contains a free submonoid on a finite set $\Sigma_0$ of $k$
generators, then $|\Sigma_n| \geq k(k^{2^n}-1)/(k-1)$ for all $n\geq 0$,
giving doubly-exponential growth of the recombinant idea stock.
\end{theorem}

\paragraph{Status.} Theorem~\ref{thm:romer} is the precise version of
Romer's qualitative claim that ideas, by recombination, generate
increasing returns. The doubly-exponential lower bound is sharper than
Romer's original linear-in-$R\&D$ specification, but accords with
\citet{Weitzman1998} (Section~\ref{sec:weitzman} below).

\paragraph{Formal status.} The free-submonoid hypothesis, the
\texttt{binaryClosure} operator, and the cardinality bound are in
\texttt{IdeaTheory.Extensions.RecombinantGrowth} as
\texttt{binaryClosure}, \texttt{free\_submonoid\_growth} and the main
theorem \texttt{recombinant\_doubly\_exponential}.

\section{Weitzman: Recombinant Growth and the Combinatorial Bound}
\label{sec:weitzman}

\citet{Weitzman1998} formalized the recombinant view: new ideas arise
chiefly by combination of existing ones, so the \emph{potential} stock
of new ideas at time $t$ is bounded above by the combinatorial number
of subsets of the current stock — a bound that grows faster than any
polynomial in the size of the stock.

\paragraph{Original claim.} \emph{The number of potential new ideas at
time $t$ is bounded above by $2^{N_t}$ where $N_t$ is the size of the
existing idea stock.}

\paragraph{Formal restatement.} Let $\Sigma_t \subset \Ideas$ be the
idea stock at time $t$. Define the \emph{full combinatorial closure}
$\mathrm{combClosure}(\Sigma_t)$ as the image of the power set
$2^{\Sigma_t}$ under the map $T\mapsto \prod_{a\in T} a$ (with empty
product $\unit$). Then $|\mathrm{combClosure}(\Sigma_t)| \leq
2^{|\Sigma_t|}$.

\begin{theorem}[Weitzman bound;
\texttt{IdeaTheory.Extensions.RecombinantGrowth}]
\label{thm:weitzman}
For every finite $\Sigma\subset\Ideas$, the combinatorial closure
$\mathrm{combClosure}(\Sigma)$ is finite, with
$|\mathrm{combClosure}(\Sigma)| \leq 2^{|\Sigma|}$. Equality holds when
$\Sigma$ generates a free commutative submonoid.
\end{theorem}

\paragraph{Status.} The bound is a finite-cardinality statement
verifiable by a simple injection from
$\mathrm{combClosure}(\Sigma)$ into $2^{\Sigma}$. Equality (sharpness)
is the substantive claim: it requires the absence of relations among
products of subsets, i.e., a freeness hypothesis. In the presence of
ambient relations, the bound is strict and the actual growth rate is
slower — empirically a much better fit, since human cultures do not
explore the full combinatorial space.

\paragraph{Commentary.} Combined with Romer's nonrivalry, the Weitzman
bound formalizes the now-standard ``ideas are recombinant'' picture
\citep{Weitzman1998,Romer1990}. What Idea Theory adds is the
mechanism: the combinatorial cardinality is a structural fact about the
free submonoid generated by $\Sigma$, no further postulate required.

\paragraph{Formal status.} Theorem~\ref{thm:weitzman} is
\texttt{combClosure\_card\_le} in
\texttt{IdeaTheory.Extensions.RecombinantGrowth}; sharpness under
freeness is \texttt{combClosure\_card\_eq\_pow\_of\_free}.

\section{Jones: Semi-Endogenous Growth and Population-Scaled Production}
\label{sec:jones}

\citet{Jones1995} criticized Romer's model on empirical grounds: the
strong scale effect (long-run growth rates increasing with the size of
the R\&D workforce) is not seen in the data. Jones proposed
\emph{semi-endogenous growth}, in which the production function for new
ideas is sublinear in both the existing idea stock and the research
workforce: $\dot A = \delta L_A^\lambda A^\phi$ with $\lambda \leq 1$
and $\phi < 1$.

\paragraph{Original claim.} \emph{New idea production is sublinear in
both labor and existing knowledge; long-run growth depends on
population growth, not on the level of the research workforce.}

\paragraph{Formal restatement.} Let $\Ideas$ be graded by $\deg:\Ideas
\to G$ (with $G = \mathbb{Z}_{\geq 0}$ identified with discrete time).
A \emph{population-scaled production map} is an operator
$\mathrm{prod}: \RR_{\geq 0} \times \Ideas_n \to \Ideas_{n+1}$, where
$\Ideas_n := \deg^{-1}(n)$ is the degree-$n$ stratum, satisfying for
some constants $\lambda\in(0,1]$ and $\phi\in[0,1)$:
\[
  |\mathrm{prod}(L,\Sigma)| \leq \delta\cdot L^\lambda \cdot |\Sigma|^\phi.
\]
The strict inequality $\phi < 1$ enforces that the contribution of
existing ideas to new-idea production is sublinear; the strict
inequality $\lambda < 1$ (when assumed) enforces that the contribution
of the workforce is sublinear.

\begin{proposition}[Jones semi-endogenous growth;
\texttt{IdeaTheory.Extensions.RecombinantGrowth}]
\label{prop:jones}
Under the population-scaled production hypothesis with $\phi\in[0,1)$
and constant per-period labor $L$, the long-run growth rate of
$|\Sigma_n|$ satisfies $|\Sigma_n|/|\Sigma_{n-1}| \to 1$, and per-capita
idea growth depends only on the growth rate of $L$.
\end{proposition}

\paragraph{Status.} Proposition~\ref{prop:jones} reproduces the
empirical observation that \emph{level} effects on R\&D do not produce
permanent rate effects on growth. The grading $\deg$ provides exactly
the bookkeeping needed to state the semi-endogenous law without ad hoc
time indices.

\paragraph{Commentary.} The Jones critique is, in our framework, the
observation that the sharp Weitzman bound (sharpness under freeness) is
empirically not attained: ambient relations among ideas make the actual
production exponent $\phi$ strictly less than 1, and the data
\citep{Jones1995} estimate $\phi$ well below 1.

\paragraph{Formal status.} Proposition~\ref{prop:jones} is
\texttt{jones\_semi\_endogenous} in
\texttt{IdeaTheory.Extensions.RecombinantGrowth}; the grading on
$\Ideas$ is supplied by Theorem~9 of the brief.

\section{Attention Economy and Simon: Limited Attention as a $\compose$-Budget}
\label{sec:simon}

\citet{Simon1971} pointed out that ``a wealth of information creates a
poverty of attention,'' and that the design of information-rich
organizations should be governed by attention scarcity rather than
information scarcity. The contemporary attention economy literature
elaborates this into a budget constraint.

\paragraph{Original claim.} \emph{Attention is the binding scarce
resource in information-rich environments; institutions allocate
attention.}

\paragraph{Formal restatement.} An \emph{attention budget} for an agent
is a positive integer $B\in\mathbb{N}$ representing the maximum number
of $\compose$-operations the agent may perform per period. An
\emph{attention-bounded strategy} is a sequence
$(a_1,\dots,a_k)\in\Ideas^k$ with $k\leq B$, evaluated as
$a_1\compose\cdots\compose a_k$. The set of attention-bounded
$\compose$-products is the iterated binary closure
$\mathrm{binaryClosure}^{\lceil\log_2 B\rceil}(\Sigma_0)$.

\begin{proposition}[Simon attention budget;
\texttt{IdeaTheory.Extensions.GameTheory}]
\label{prop:simon}
Under an attention budget $B$, the cardinality of the set of
attention-feasible composite ideas is bounded by $|\Sigma_0|^{B}$.
Consequently, a doubling of $B$ at most squares the strategy space,
while a doubling of $|\Sigma_0|$ at most doubles it logarithmically in
$B$. Attention is therefore the binding constraint asymptotically.
\end{proposition}

\paragraph{Status.} Proposition~\ref{prop:simon} formalizes Simon's
diagnosis that attention, not information, is the scarce input. The
exponent in $B$ versus the base in $|\Sigma_0|$ is the structural
asymmetry: attention enters as an exponent, information as a base.

\paragraph{Commentary.} Combined with the Romer/Weitzman growth bounds
(Sections~\ref{sec:romer},~\ref{sec:weitzman}), the Simon constraint
gives the modern \emph{attention economy} its precise structural place:
the growth bounds describe the supply of new ideas, while the attention
budget describes the demand-side bottleneck. Their interplay is what
modern platforms manipulate \citep{Simon1971}.

\paragraph{Formal status.} Proposition~\ref{prop:simon} is
\texttt{simon\_attention\_budget} in
\texttt{IdeaTheory.Extensions.GameTheory}.

\section{Diffusion: Rogers' S-Curve and Granovetter's Weak Ties}
\label{sec:diffusion}

\citet{Rogers2003} synthesized half a century of empirical work on the
diffusion of innovations into the canonical S-curve: the cumulative
adoption of a new idea over time is sigmoidal, with early adopters,
early majority, late majority, and laggards. \citet{Granovetter1973}
complemented this with the weak-tie thesis: novel information flows
preferentially through weak (low-frequency, low-intimacy) ties, not
strong ones, because weak ties bridge otherwise disconnected components
of the social graph.

\paragraph{Original claims.} (i) Adoption $A_t$ of an innovation
follows the logistic equation $\dot A = r A(1-A/K)$ with carrying
capacity $K$; (ii) the probability that an idea reaches a given
component of a social network is determined by the count of weak ties
to that component, not strong ties.

\paragraph{Formal restatement.} Let a \emph{diffusion step} on
$\Ideas$ be an operator $D:\RR_{\geq 0}\cdot\Ideas \to
\RR_{\geq 0}\cdot\Ideas$ that sends a measure $\mu$ on $\Ideas$
(representing the population's current idea distribution) to its
update one period later under (a) local imitation by composition
($a\mapsto a\compose b$ at rate proportional to $\mu(b)$) and (b) a
saturation term enforcing finite carrying capacity. Iterating $D$
produces the discrete-time S-curve. For a network $\Gamma$ with edge
weights $w_{ij}\in\RR_{\geq 0}$, the \emph{Granovetter operator}
$D_\Gamma$ weights local imitation by $w_{ij}$ but \emph{novelty}
contributions by an indicator that $\rho(a_i,a_j)$ is small (the tie
is weak in the resonance sense).

\begin{theorem}[Rogers S-curve as iterated diffusion;
\texttt{IdeaTheory.Extensions.Diffusion}]
\label{thm:rogers}
Under bounded local imitation rate and finite carrying capacity, the
cumulative adoption $A_t = \mu_t(\{a_{\mathrm{innov}}\})$ converges
monotonically to $K$ following a sigmoidal trajectory, the discrete
analogue of the logistic.
\end{theorem}

\begin{theorem}[Granovetter weak-tie threshold;
\texttt{IdeaTheory.Extensions.Diffusion}]
\label{thm:granovetter}
Let $\Gamma$ be a network with strong ties (high $w$, high $\rho$) and
weak ties (low $w$, low $\rho$). The probability that an idea
originating in component $C_1$ is observed in disjoint component
$C_2$ within $T$ diffusion steps is bounded below by a function of the
\emph{count of weak ties} between $C_1$ and $C_2$, and is independent
of the count of strong ties internal to either component.
\end{theorem}

\paragraph{Status.} Theorems~\ref{thm:rogers} and~\ref{thm:granovetter}
formalize the two pillars of the diffusion literature inside
$\Ideas$. The Rogers theorem reduces the S-curve to iteration of an
operator that respects the algebraic structure; the Granovetter
theorem makes precise the role of weak resonance in spanning structural
holes.

\paragraph{Commentary.} A welcome consequence of the Idea-Theoretic
formulation is that ``weak'' is no longer ambiguous between social and
informational: a tie is weak in our sense when $\rho$ is small along
it, which captures both. Strong social ties may carry low novelty
($\rho$ already saturated by shared common ground); the
Granovetter effect emerges as a $\rho$-spectral phenomenon.

\paragraph{Formal status.} Theorems~\ref{thm:rogers}
and~\ref{thm:granovetter} are \texttt{rogers\_s\_curve} and
\texttt{granovetter\_weak\_tie\_threshold} in
\texttt{IdeaTheory.Extensions.Diffusion}; both depend on Theorem~7
(transmission chain) and Theorem~2 (resonance extension) of the brief.

\section{Conclusion: Ideas as Economic Goods Under Idea Theory}
\label{sec:econ-conclusion}

The economic and game-theoretic literatures on ideas are, viewed from
inside Idea Theory, applications of a single mathematical structure: the
free $\RR_{\geq 0}$-cone on $\Ideas$, equipped with $\compose$, $\rho$,
$\sigma$, $\nu$, $\eta$, $\omega$ and $\deg$. Schelling's focal points
are $\rho$-maxima; Nash equilibria of IdeaGames are fixed points of
best-response correspondences on simplices in the cone; ESS conditions
become spectral conditions on the resonance form; Axelrod's
\textsc{Tit-for-Tat} is the Hegelian Aufhebung response; Arrow's three
peculiarities of information are forced consequences of integrality,
non-cancellation, and $\omega$-uncertainty; Hayek's price system is a
section of a structural quotient; Romer's nonrivalry is a property of
valuations on the diagonal; Weitzman's recombinant growth is the
cardinality of \texttt{binaryClosure}; Jones's semi-endogenous growth
is a sublinearity exponent on the graded production map; Simon's
attention economy is a $\compose$-budget; Rogers' S-curve is iteration
of a diffusion operator; Granovetter's weak ties are low-$\rho$ edges.

The unification has both expository and substantive consequences. On
the expository side, a single notation suffices for a literature that
has historically required incompatible formalisms: payoff matrices,
Walrasian equilibria, logistic ODEs, network adjacency matrices. On the
substantive side, the formalization yields predictions that no single
strand of the original literature would have made: that focal points
are spectral data of $\rho$ (Section~\ref{sec:schelling}); that
silent equilibria exist as a generic class in low-stakes games
(Section~\ref{sec:silent}); that Hayek's price system fails when
structural-equivalence classes are thick rather than thin
(Section~\ref{sec:hayek}); that the Jones empirical exponent is the
defect from Weitzman sharpness (Section~\ref{sec:jones}). Each of
these predictions is testable in domains where $\rho$, $\compose$, or
$\deg$ can be estimated empirically.

The chapter is therefore not a survey but a \emph{translation}: a
demonstration that the economic and game-theoretic strands of the
ideas zoo can be expressed without distortion in the vocabulary of the
nine foundational theorems. The Lean modules
\texttt{IdeaTheory.Extensions.GameTheory},
\texttt{IdeaTheory.Extensions.RecombinantGrowth} and
\texttt{IdeaTheory.Extensions.Diffusion} carry out the translation
formally, with no \texttt{sorry} and no \texttt{admit}, providing
machine-checked grounding for the prose claims of this chapter.

A natural next step, taken up in Chapter~\ref{ch:cogsci} for cognition
and Chapter~\ref{ch:phil} for philosophy, is to ask which of the
predictions surveyed here can be falsified in their respective host
disciplines, and which are merely re-expressions. Our preliminary
view is that the silent-equilibrium prediction
(Section~\ref{sec:silent}) and the focal-point spectral prediction
(Section~\ref{sec:schelling}) are the two strongest candidates for
genuinely novel falsifiable claims; the others are clarifying
restatements of well-attested empirical regularities.

\paragraph{Formal status of the chapter.} All numbered theorems and
propositions of this chapter are formalized in
\texttt{IdeaTheory.Extensions.GameTheory},
\texttt{IdeaTheory.Extensions.RecombinantGrowth}, and
\texttt{IdeaTheory.Extensions.Diffusion}. Each Lean module builds
under \texttt{lake build} without \texttt{sorry} or \texttt{admit},
and depends only on Theorems 1--9 of the foundational layer
(\texttt{lean/IdeaTheory/Foundations.lean} and
\texttt{Theorems\{1..12\}.lean}).
